\section{The direct realization and its analytical purpose}\label{sec:intro}

The direct problem is to realize the Newman deformation through a specified
preparation, evolution and signed observation, and then to understand what
that realization says about zero collisions. The shared manuscript
\cite{TNDBackground} supplies the geometric and spectral background. Here we
retain only the background needed for this branch and derive its preparation,
complete instrument, calibrated observation and collision diagnostics in one
continuous development.

The canonical family is
\begin{align}
 \Phi(v)&=\sum_{n\ge1}(2\pi^2n^4\ee^{9v}-3\pi n^2\ee^{5v})
                       \ee^{-\pi n^2\ee^{4v}},\label{eq:Phi}\\
 H_\lambda(z)&=\int_0^\infty\ee^{\lambda v^2}\Phi(v)\cos(zv)\dd v.
 \label{eq:H}
\end{align}
The complete source is smooth, positive and even; its reflected interpretation
and all limits are justified below. We keep the classical normalization
\begin{equation}\label{eq:xi}
 H_0(z)=\frac18\xi\!\left(\frac{1+\ii z}{2}\right),\qquad
 \xi(s)=\frac12s(s-1)\pi^{-s/2}\Gamma(s/2)\zeta(s),
\end{equation}
as in \cite[Introduction, (1)--(4)]{RodgersTao2020}. No theorem about the
zeros of $H_0$ is used. Direct differentiation gives
$\partial_\lambda H_\lambda=-\partial_z^2H_\lambda$.

The model keeps the free kinetic Hamiltonian on
$\mathcal C=\R\times S^2_{1/2}$. It adds one prescribed count instrument
that monitors the freely transported position quadrature. Its no-count
branch attenuates the positive square-root source by
$\ee^{-\gamma tX^2/2}$ in the interaction frame. The resulting density has
parameter $\lambda(t)=1/2-\gamma t$ on $0\le t\le1/\gamma$.
Physical time therefore runs toward decreasing Newman parameter.

The principal identities organize the paper. If $p_t$ is the no-count
probability, $C(t,k)$ the conditional cosine expectation and $q_\pm$ the joint
success/outcome probabilities, then
\begin{equation}\label{eq:mainintro}
 H_{\lambda(t)}(k)=H_{1/2}(0)\,p_t C(t,k)
                 =H_{1/2}(0)[q_+(t,k)-q_-(t,k)].
\end{equation}
The full instrument has outputs $\sigma_j(t)$, one for each count $j$.
Their signed traces $d_j(t,k)=\Tr[\sigma_j(t)\cos(kQ_t)]$ satisfy
\begin{equation}\label{eq:countintro}
 d_j=\frac{(-\gamma t)^j}{j!}\partial_k^{2j}d_0,\qquad
 \partial_td_0=-d_1/t\quad(t>0),\qquad
 \sum_{j\ge0}d_j=C(0,k).
\end{equation}
Thus the complete count record organizes weighted source moments, while the
Newman evolution appears in its selected no-count sector.

At a hypothetical nondegenerate double event,
$d_0=\partial_kd_0=0$ and $a_*=\partial_k^2d_0\ne0$. The local law is
\begin{equation}\label{eq:foldintro}
 d_0(t_*+h,k_*+x)=\frac{a_*}{2}(x^2+2\gamma h)
                      +O(|x|^3+|hx|+h^2).
\end{equation}
The positive state and success mass remain regular. Equal no-count outcome
probabilities have unequal rates of decrease, encoded by
$d_1(t_*,k_*)=-\gamma t_*a_*$. This provides a dynamical description of a
collision, without asserting that the canonical family has one.

\subsection{Assumptions, evidence and organization}
The source loading, its identification with physical position, the complete
angular preparation, the initial scale $H_{1/2}(0)$ and access to the
specified no-count cosine measurement are supplied model data. The proofs
establish the consistency and exact observation of that model. They do not
derive a microscopic apparatus or an intrinsic theta loader.

The complete realization \cite{TNDDirect} has an internal assessment
accepting its conditional scope \cite{TNDDirectReview}; the shared background
has its own internal assessment \cite{TNDSharedReview}. The later count and
collision analysis \cite{TNDDirectCollision} retains proof-author diagnostic
status. The exposition below makes its derivations explicit but is not an
independent reassessment. No numerical evidence is needed for any proof.

Section~\ref{sec:preparation} establishes the domains and admissible prepared
state. Sections~\ref{sec:instrument}--\ref{sec:observation} construct the
instrument and prove the absolute readout. Section~\ref{sec:count} derives
the complete count calculus. Sections~\ref{sec:collisions}--\ref{sec:control}
analyze local collisions and an exact colliding control. The final section
states the outstanding tasks. The distinction between a representation and
a collision-preventing mechanism remains visible throughout.
